\begingroup
\renewcommand{\arraystretch}{1.45}
\Needspace{45\baselineskip}
\InterpolationArticle{28}
\textit{Beispiel.} In der vom angesehenen Baron \textsc{de Zach} herausgegebenen gelehrten Korrespondenz, Band X, S. 188, findet sich eine Tafel, welche die nördliche und südliche Grenze des Zodiaks der \textit{Pallas} angibt. Die Deklination jeder der beiden Grenzen wird als periodische Funktion der Rektaszension betrachtet, die in jener Tafel in Schritten von fünf Grad fortschreitet. Zur Veranschaulichung der vorangehenden Untersuchungen werden wir hier eine Anwendung auf die nördliche Grenze machen; deren Deklination wird also für uns $X$ und deren Rektaszension $x$ sein. Aus jener Tafel entnehmen wir die folgende Periode von 12 Gliedern:
{\small\setlength{\arraycolsep}{3pt}\[
\begin{array}{r|r r l c r}
x&\multicolumn{5}{c}{X}\\ \hline
0^{\mathord{0}}&6^{\mathord{0}}&48'&\text{N.}&=+&408'\\
30&1&29&.\,.\,.\,.\,.\,.\,.&+&89\\
60&1&6&\text{S.}\,.\,.&-&66\\
90&0&10&\text{N.}\,.\,.\,.&+&10\\
120&5&38&.\,.\,.\,.\,.\,.\,.&+&338\\
150&13&27&.\,.\,.\,.\,.\,.\,.&+&807\\
180&20&38&.\,.\,.\,.\,.\,.\,.&+&1238\\
210&25&11&.\,.\,.\,.\,.\,.\,.&+&1511\\
240&26&23&.\,.\,.\,.\,.\,.\,.&+&1583\\
270&24&22&.\,.\,.\,.\,.\,.\,.&+&1462\\
300&19&43&.\,.\,.\,.\,.\,.\,.&+&1183\\
330&13&24&.\,.\,.\,.\,.\,.\,.&+&804
\end{array}
\]}
\SourceLeaf{318}
\Needspace{47\baselineskip}
Teilen wir diese Periode zunächst in drei Perioden von je vier Gliedern:
{\small\setlength{\arraycolsep}{3pt}\[
\begin{array}{r r|r r|r r}
a=0^{\mathord{0}}&A=+408&a'=30^{\mathord{0}}&A'=+89&a''=60^{\mathord{0}}&A''=-66\\
b=90^{\mathord{0}}&B=+10&b'=120^{\mathord{0}}&B'=+338&b''=150^{\mathord{0}}&B''=+807\\
c=180^{\mathord{0}}&C=+1238&c'=210^{\mathord{0}}&C'=+1511&c''=240^{\mathord{0}}&C''=+1583\\
d=270^{\mathord{0}}&D=+1462&d'=300^{\mathord{0}}&D'=+1183&d''=330^{\mathord{0}}&D''=+804
\end{array}
\]}
In der Formel
\[
\begin{aligned}
&X'={}\gamma+\gamma'\cos x+\gamma''\cos2x\\[5pt]
&+\delta'\sin x+\delta''\sin2x
\end{aligned}
\]
ergibt sich also
\[
\begin{aligned}
&\gamma=\tfrac14(A+B+C+D)=779'{,}5\\[5pt]
&\gamma'=\tfrac12(A\cos a+B\cos b+C\cos c+D\cos d)\\[5pt]
&=\tfrac12(A-C)=-415'{,}0\\[5pt]
&\delta'=\tfrac12(A\sin a+B\sin b+C\sin c+D\sin d)\\[5pt]
&=\tfrac12(B-D)=-726'{,}0\\[5pt]
&\gamma''=\tfrac14(A\cos2a+B\cos2b+C\cos2c+D\cos2d)\\[5pt]
&=\tfrac14(A-B+C-D)\\[5pt]
&=+43'{,}5\\[5pt]
&\delta''=\tfrac14(A\sin2a+B\sin2b+C\sin2c+D\sin2d)=0
\end{aligned}
\]
und entsprechend für die zweite und dritte Periode. Auf diese Weise entsteht
{\small\setlength{\arraycolsep}{3pt}\[
\begin{array}{l|r|r|r|r|r|r}
\text{Für die Periode}&\text{bei }y=4x&\gamma&\gamma'&\delta'&\gamma''&\delta''\\ \hline
\text{erste}&0^{\mathord{0}}&+779{,}5&-415{,}0&-726{,}0&+43{,}5&0\\
\text{zweite}&120^{\mathord{0}}&+780{,}2&-404{,}5&-721{,}4&+9{,}9&+17{,}1\\
\text{dritte}&240^{\mathord{0}}&+782{,}0&-413{,}5&-713{,}3&+11{,}7&-20{,}3
\end{array}
\]}
\Needspace{42\baselineskip}
Hier wird ferner $\gamma$ durch die Formel
\[
\begin{aligned}
&\tfrac13(779'{,}5+780{,}2+782{,}0)\\[5pt]
&+\tfrac23(779{,}5+780{,}2\cos120^{\mathord{0}}+782{,}0\cos240^{\mathord{0}})\cos4x\\[5pt]
&+\tfrac23(780{,}2\sin120^{\mathord{0}}+782{,}0\sin240^{\mathord{0}})\sin4x
\end{aligned}
\]
oder durch
\[
\begin{aligned}
&780{,}6-1{,}1\cos4x-1{,}0\sin4x,\quad\text{und ebenso}
\end{aligned}
\]
\[
\begin{aligned}
&\gamma'\quad\text{durch}\quad-411{,}0-4{,}0\cos4x+5{,}2\sin4x\\[5pt]
&\delta'\quad\text{durch}\quad-720{,}2-5{,}8\cos4x-4{,}7\sin4x\\[5pt]
&\gamma''\quad\text{durch}\quad+21{,}7+21{,}8\cos4x-1{,}1\sin4x\\[5pt]
&\delta''\quad\text{durch}\quad-1{,}1+1{,}1\cos4x+21{,}6\sin4x
\end{aligned}
\]
dargestellt. Werden diese Werte in $X'$ eingesetzt, so entsteht die folgende Formel, welche die 12 vorgegebenen Werte darstellt:
\SourceLeaf{319}
\[
\begin{aligned}
&+780{,}6\\[5pt]
&-411{,}0\cos x-720{,}2\sin x\\[5pt]
&+43{,}4\cos2x-2{,}2\sin2x\\[5pt]
&-4{,}3\cos3x+5{,}5\sin3x\\[5pt]
&-1{,}1\cos4x-1.0\sin4x\\[5pt]
&+0{,}3\cos5x-0{,}3\sin5x\\[5pt]
&+0{,}1\cos6x
\end{aligned}
\]
\Needspace{42\baselineskip}
Teilen wir zweitens dieselbe Periode in vier Perioden von je drei Gliedern, deren jeder also durch eine Formel erster Ordnung
\[
\begin{aligned}
&X'=\gamma+\gamma'\cos x+\delta'\sin x
\end{aligned}
\]
genügt werden muss. Hier findet sich
\begin{center}
\begin{tabular}{r|c|r|r|r}
Für die Periode & bei $y=3x$ & $\gamma$ & $\gamma'$ & $\delta'$\\\hline
erste & $0^{0}$ & $+776{,}3$ & $-368{,}3$ & $-718{,}8$\\
zweite & $90^{0}$ & $+786{,}0$ & $-414{,}5$ & $-676{,}0$\\
dritte & $180^{0}$ & $+785{,}0$ & $-453{,}0$ & $-721{,}1$\\
vierte & $270^{0}$ & $+775{,}0$ & $-408{,}2$ & $-765{,}0$
\end{tabular}
\end{center}
woraus sich die folgenden Formeln zur Darstellung von $\gamma,\gamma',\delta'$ herleiten:
\[
\begin{aligned}
&\gamma\quad.\,.\,.\,.\quad+780{,}6-4{,}3\cos3x+5{,}5\sin3x\\[5pt]
&\qquad{}+0{,}1\cos6x\\[5pt]
&\gamma'\quad.\,.\,.\,.\quad-411{,}0+42{,}3\cos3x-3{,}2\sin3x\\[5pt]
&\qquad{}+0{,}3\cos6x\\[5pt]
&\delta'\quad.\,.\,.\,.\quad-720{,}2+1{,}2\cos3x+44{,}5\sin3x\\[5pt]
&\qquad{}+0{,}3\cos6x
\end{aligned}
\]
Werden diese Werte in $\gamma+\gamma'\cos x+\delta'\sin x$ eingesetzt und die letzten Teile in $\gamma,\gamma',\delta'$, die $\cos6x$ enthalten, nach den Vorschriften von Artikel 26 behandelt, so ergibt sich genau dieselbe Formel, zu der wir oben gelangt sind.
\endgroup
